\section*{Chapter \ref{ch:b2:intertemporal-choice} --- Intertemporal Choice}

\begin{solution}{exo:b2:intertemporal-choice:1}
% code: exercises/solutions.py ex_pv
$\num{30000} + \num{40000}/1.04 + \num{50000}/1.04^2 = \num{30000} + \num{38461.5} + \num{46227.8} =
\dal{114689}$.
\end{solution}

\begin{solution}{exo:b2:intertemporal-choice:2}
% code: exercises/solutions.py ex_euler_log
Lifetime resources $\num{40000} + \num{42000}/1.05 = \num{80000}$. With $u = \ln$ the \omterm{def:b2:intertemporal-choice:euler-equation}{Euler equation} is
$c_1 = \beta(1 + r)c_0 = 0.9975\,c_0$, and the budget $c_0 + c_1/1.05 = c_0(1 + \beta) = \num{80000}$.
So $c_0 = \num{80000}/1.95 = \dal{41026}$, $c_1 = \dal{40923}$, and saving is $\num{40000} -
\num{41026} = -\dal{1026}$: the household borrows.
\end{solution}

\begin{solution}{exo:b2:intertemporal-choice:3}
% figdata: ce_age (ratio, spending_pp)
(a) Spending over income after taxes: $61.59/84.94 = 0.725$ at 25--34, and $51.15/49.98 = 1.023$ at 75
and older. (b) $44.47/2.0 = 22.24$ thousand dollars per person for the youngest group, $72.25/2.2 = 32.84$
for the 55--64 group. (c) Per person: spending per person rises 1.48 times from the youngest group to
the 55--64 group, while spending per unit rises 1.85 times to its peak and income after taxes 2.26
times. Much of the hump in spending per unit is the number of people in the unit.
\end{solution}

\begin{solution}{exo:b2:intertemporal-choice:4}
Refutations: a preference reversal (choosing the larger, later reward when both are distant and the
smaller, sooner one when it is immediate, with the same delay between them); paying a positive price
for a \omterm{def:b2:intertemporal-choice:commitment-device}{commitment device} (\cref{prop:b2:intertemporal-choice:commitment}); a plan abandoned with no
new information. Falling consumption over life is what the model predicts when $\beta(1 + r) < 1$ (an
impatient household, \cref{cor:b2:intertemporal-choice:crra-growth}), so it refutes nothing.
\end{solution}

\begin{solution}{exo:b2:intertemporal-choice:5}
% code: exercises/solutions.py ex_crra_growth
$\gamma = 2$: approximately $(0.04 - 0.02)/2 = 1$ per cent; exactly $(1.04/1.02)^{1/2} - 1 = 0.976$
per cent. $\gamma = 0.5$: approximately $0.02/0.5 = 4$ per cent; exactly $(1.04/1.02)^{2} - 1 = 3.96$
per cent. The second household has the larger EIS, $1/\gamma = 2$ against $0.5$: its consumption
responds four times as much to the gap between $r$ and $\rho$.
\end{solution}

\begin{solution}{exo:b2:intertemporal-choice:6}
% code: exercises/solutions.py ex_rate_rise
By \cref{prop:b2:intertemporal-choice:rate}, saving rises for $\gamma = 0.5$, is unchanged for $\gamma =
1$ and falls for $\gamma = 2$. With $c_0 = \num{100000}/(1 + 0.96^{1/\gamma}(1 + r)^{1/\gamma - 1})$:
for $\gamma = 2$, $c_0$ goes from \dal{50880} to \dal{51238} (the income effect wins); for $\gamma = 1$
it stays at $\num{100000}/1.96 = \dal{51020}$ (the two effects cancel); for $\gamma = 0.5$ it falls from
\dal{51302} to \dal{50584} (the substitution effect wins). There is no human-wealth effect, since all
income comes in period 0.
\end{solution}

\begin{solution}{exo:b2:intertemporal-choice:7}
% code: exercises/solutions.py ex_borrowing
(a) Resources $\num{20000} + \num{60000}/1.05 = \num{77143}$; $c_0 = \num{77143}/1.95 = \dal{39560}$
and $c_1 = 0.9975\,c_0 = \dal{39462}$, so the student borrows \dal{19560}.
(b) At the endowment, $u'(c_0)/u'(c_1) = \num{60000}/\num{20000} = 3$, far above $\beta(1 + r) = 0.9975$: the Euler
inequality holds strictly, the student wants to move consumption towards the present, and the limit
binds. (c) Its income: \dal{20000} now and \dal{60000} next year.
\end{solution}

\begin{solution}{exo:b2:intertemporal-choice:8}
% code: exercises/solutions.py ex_brenland
Band means, disposable income and consumption in dalens, and their ratio: 21--29: \num{21078},
\num{20033}, $0.950$; 30--39: \num{25501}, \num{23852}, $0.935$; 40--49: \num{28378}, \num{26371},
$0.929$; 50--59: \num{29813}, \num{27536}, $0.924$; 60--69: \num{24926}, \num{23422}, $0.940$;
70--79: \num{18654}, \num{18143}, $0.973$. The least-squares slope of consumption on disposable
income across households is $0.844$. Consumption tracks income closely: the highest band mean is 1.60
times the lowest for income and 1.52 times for consumption. Smoothing predicts a much flatter
consumption profile than income profile (flat if $\beta(1 + r) = 1$), with the old spending more than
their income and the middle-aged much less. Brenland's households look like households at their
borrowing limit, or households that do not plan.
\end{solution}

\begin{solution}{exo:b2:intertemporal-choice:9}
(a) The household's later selves carry out its date-0 plan with or without the account
(\cref{prop:b2:intertemporal-choice:consistent}); the account removes options and costs \dal{20}, so
it can only lower utility. (b) The naive household believes its later selves will follow its plan, so
it sees no reason to restrict them and does not pay the fee, although it would gain from the account.
(c) Augenblick, Niederle and Sprenger: 59 per cent took a \omterm{def:b2:intertemporal-choice:commitment-device}{commitment device} when it was free; almost
nobody paid 25 cents for it. The gym members did not buy commitment: they paid for flexibility they
did not use.
\end{solution}

\begin{solution}{exo:b2:intertemporal-choice:10}
Seen from date 0, the marginal rate of substitution between dates $t + k$ and $t + k + 1$ is
$D(t + k)u'(c_{t+k})/(D(t + k + 1)u'(c_{t+k+1}))$; seen from date $t$ it is
$D(k)u'(c_{t+k})/(D(k + 1)u'(c_{t+k+1}))$. A date-0 optimal plan sets the first equal to $1 + r$ (by
the Euler argument); for the date-$t$ self to keep the plan for every $u$ and every budget, the second
must equal $1 + r$ at the same consumptions. So $D(t + k + 1)/D(t + k) = D(k + 1)/D(k)$ for all $t$ and
$k$. With $k = 0$: $D(t + 1)/D(t) = D(1)/D(0) = D(1)$ for every $t$, so $D(t) = D(1)^{t}$. Exponential
discounting satisfies the condition (\cref{prop:b2:intertemporal-choice:consistent}), and nothing else
does.
\end{solution}

\begin{solution}{exo:b2:intertemporal-choice:11}
% code: exercises/solutions.py ex_cake
(a) The date-1 self maximises $\ln c_1 + 0.5\ln c_2$ with $c_1 + c_2 = W_1$: $c_1 = W_1/1.5$.
(b) The date-0 self then values $\ln c_0 + 0.5(\ln c_1 + \ln c_2) = \ln c_0 + 0.5 \times 2\ln W_1 +$ a
constant, with $W_1 = \num{90000} - c_0$: $1/c_0 = 1/(\num{90000} - c_0)$, so $c_0 = \num{90000}/(1 +
2\varphi) = \dal{45000}$, $W_1 = \num{45000}$, $c_1 = \dal{30000}$, $c_2 = \dal{15000}$. (c) The date-0
self weights periods 1 and 2 equally, so it wants $c_1 = c_2 = \dal{22500}$. (d) A naive date-0 self
expects $c_1 = c_2 = W_1/2$, and values $\ln c_0 + 2 \times 0.5\ln(W_1/2)$: the same function of $c_0$
up to a constant, so the same $c_0 = \dal{45000}$. With logarithmic utility, what the date-0 self
believes about later choices changes the constant, not the first decision.
\end{solution}

\begin{solution}{exo:b2:intertemporal-choice:12}
% code: exercises/solutions.py ex_borrowing
Unconstrained utility $\ln\num{39560} + 0.95\ln\num{39462}$; constrained utility $\ln\num{20000} +
0.95\ln\num{60000}$. Scaling resources by $\lambda$ scales both consumptions by $\lambda$ and adds
$(1 + \beta)\ln\lambda$ to utility, so $\ln\lambda = (U_{\text{limit}} - U_{\text{free}})/1.95$, which
gives $\lambda = 0.86445$. The limit costs the student as much as losing 13.6 per cent of its lifetime
resources: $0.13555 \times \num{77143} = \dal{10456}$.
\end{solution}

\begin{solution}{pb:b2:intertemporal-choice:1}
% code: exercises/solutions.py problem
\textbf{1.} $b_0 = m_0 - c_0$, $b_1 = 1.25\,b_0 + m_1 - c_1$, $0 = b_2 = 1.25\,b_1 + m_2 - c_2$.
Dividing the last by $1.25^2$ and the middle one by $1.25$ and adding: $c_0 + c_1/1.25 + c_2/1.25^2 =
m_0 + m_1/1.25 + m_2/1.25^2$.

\textbf{2.} $\mathcal{M} = \num{24000} + \num{60000}/1.25 + \num{6000}/1.5625 = \num{24000} +
\num{48000} + \num{3840} = \dal{75840}$.

\textbf{3.} $\rho = 1/\beta - 1 = 25$ per cent, equal to $r$: $\beta(1 + r) = 1$.

\textbf{4.} $\ln\num{24000} + 0.8\ln\num{60000} + 0.64\ln\num{6000} = 24.4552$. The weights add up to
$2.44$, so $\bar c_h = \exp(24.4552/2.44) = \dal{22530}$.

\textbf{5.} $1/c_0 = 0.8 \times 1.25/c_1$ and $1/c_1 = 0.8 \times 1.25/c_2$: $c_0 = c_1 = c_2$.

\textbf{6.} $c(1 + 0.8 + 0.64) = \num{75840}$, so $c = \dal{31082}$.

\textbf{7.} $b_0 = \num{24000} - \num{31082} = -\dal{7082}$; $b_1 = 1.25 \times (-\num{7082}) +
\num{60000} - \num{31082} = \dal{20066}$; $1.25 \times \num{20066} + \num{6000} - \num{31082} = 0$.

\textbf{8.} $\num{31082} - \num{22530} = \dal{8552}$ per period.

\textbf{9.} Growth factor $\beta(1 + r) = 0.8 \times 1.5 = 1.2$. $\mathcal{M} = \num{24000} +
\num{60000}/1.5 + \num{6000}/2.25 = \dal{66667}$. $c_0 = \num{66667}/2.44 = \dal{27322}$, $c_1 =
\dal{32787}$, $c_2 = \dal{39344}$.

\textbf{10.} $b_0 = \num{24000} - \num{27322} = -\dal{3322}$: she borrows less. In general the budget is
$c_0(1 + \beta(1 + r)/(1 + r) + \beta^2(1 + r)^2/(1 + r)^2) = c_0(1 + \beta + \beta^2) = \mathcal{M}$.
The substitution and income effects cancel under $u = \ln$; what remains is the human-wealth effect:
her income comes mostly later, and the higher rate lowers its present value from \dal{75840} to
\dal{66667}.

\textbf{11.} No. With $\beta(1 + r) = 1$ the \omterm{def:b2:intertemporal-choice:euler-equation}{Euler equation} is $u'(c_t) = u'(c_{t+1})$ for any
strictly concave $u$ (\cref{prop:b2:intertemporal-choice:flat}).

\textbf{12.} Yes: the unconstrained plan has $b_0 = -\num{7082} < 0$.

\textbf{13.} $c_0 = \dal{24000}$. At date 1 she has $\num{60000} + \num{6000}/1.25 = \num{64800}$ in
present value, and $c_1 = c_2 = \num{64800}/1.8 = \dal{36000}$.

\textbf{14.} $u'(c_0) = 1/\num{24000} > \beta(1 + r)u'(c_1) = 1/\num{36000}$: strict inequality, the
limit binds. Between 1 and 2: $1/\num{36000} = 1/\num{36000}$.

\textbf{15.} With the limit $\ln\num{24000} + 1.44\ln\num{36000} = 25.1932$; without, $2.44\ln\num{31082}
= 25.2403$.

\textbf{16.} $\bar c = \exp(25.1932/2.44) = \dal{30488}$; the limit costs $\num{31082} - \num{30488} =
\dal{594}$ per period ($593.6$ before rounding), or $593.6 \times 2.44 = \dal{1448}$ of lifetime
resources (1.9 per cent of $\mathcal{M}$).

\textbf{17.} Yes: paying \dal{1000} leaves \dal{74840} and a flat plan of $\num{74840}/2.44 =
\dal{30672}$, more than the constrained \dal{30488}. She would pay up to \dal{1448}.

\textbf{18.} Date 0: $1$, $\varphi\beta = 0.56$, $\varphi\beta^2 = 0.448$. Date 1: $1$ on period 1 and
$\varphi\beta = 0.56$ on period 2.

\textbf{19.} Its weights on periods 1 and 2 are in the ratio $\beta$, and $\beta(1 + r) = 1$: $c_1 =
c_2$, so $c_1 = W_1/1.8$.

\textbf{20.} $1/c_1 = 0.56 \times 1.25/c_2$: $c_2 = 0.7\,c_1$, and $c_1 + c_2/1.25 = W_1$ gives $c_1 =
W_1/1.56$.

\textbf{21.} $c_0 = \mathcal{M}/(1 + 0.56 \times 1.8) = \num{75840}/2.008 = \dal{37769}$; $W_1 =
(\num{75840} - \num{37769}) \times 1.25 = \dal{47589}$; $c_1 = \num{47589}/1.56 = \dal{30506}$; $c_2 =
(\num{47589} - \num{30506}) \times 1.25 = \dal{21354}$.

\textbf{22.} No. A naive self expects $c_1 = c_2 = W_1/1.8$, but either way $W_1$ enters its utility as
$\varphi\beta(1 + \beta)\ln W_1$ plus a constant, so its choice of $c_0$ is the same; and the date-1
self then makes the same choice. With logarithmic utility the two types follow the same path.

\textbf{23.} She fixes $c_1 = c_2 = W_1/1.8 = \dal{26438}$, the date-0 self's preferred split. The
date-0 self values $W_1$ through $1.008\ln W_1$, so the account is worth as much to it as multiplying
$W_1$ by the $x$ for which $1.008\ln x$ equals its gain in utility, $0.56\ln\num{26438} + 0.448\ln\num{26438}
- 0.56\ln\num{30506} - 0.448\ln\num{21354}$. Computed with the unrounded consumptions, that gives $x = 1.01554$: $0.01554 \times \num{47589} =
\dal{740}$ at date 1, \dal{592} at date 0 (the rounded consumptions printed here give \dal{739} and \dal{591}).

\textbf{24.} Exponential plan (\dal{31082} each period): $25.2403$. Present-biased without the account:
$\ln\num{37769} + 0.8\ln\num{30506} + 0.64\ln\num{21354} = 25.1799$. With the account: $25.2021$. The
account fixes the split between middle age and retirement, but not the overconsumption of youth
(\dal{37769} against \dal{31082}), which it cannot reach because the young self chooses it.

\textbf{25.} A ban on borrowing when young costs the engineer \dal{594} of consumption per period,
\dal{1448} of lifetime resources; if she is present-biased, a commitment savings account opened when
young is worth about \dal{740} to her, at date 1.
\end{solution}
